% ECONOMICS_PROBLEM: {"id": "C73-1", "title": "Equilibrium payoffs in finite multiplayer stochastic games", "jel_code": "C73", "source_id": "OP-0590"}
% FIRST_POSTED: 2026-10-09
% ATTEMPT_STATUS: unresolved
% ENTRY_KIND: research_attempt
\documentclass[11pt,a4paper]{article}
\usepackage[T1]{fontenc}
\usepackage[utf8]{inputenc}
\usepackage{lmodern}
\usepackage{microtype}
\usepackage[margin=1.08in]{geometry}
\usepackage{amsmath,amssymb,amsthm,mathtools}
\usepackage{enumitem}
\usepackage[colorlinks=true,linkcolor=blue!50!black,citecolor=blue!50!black,urlcolor=blue!50!black]{hyperref}
\usepackage{xcolor}
\usepackage{booktabs}
\usepackage{array}
\usepackage{fancyhdr}
\usepackage{authblk}
\usepackage{mathrsfs}
\usepackage{bbm}
\numberwithin{equation}{section}
\newtheorem{theorem}{Theorem}[section]
\newtheorem{proposition}[theorem]{Proposition}
\newtheorem{lemma}[theorem]{Lemma}
\newtheorem{corollary}[theorem]{Corollary}
\theoremstyle{definition}
\newtheorem{definition}[theorem]{Definition}
\newtheorem{question}[theorem]{Open question}
\theoremstyle{remark}
\newtheorem{remark}[theorem]{Remark}
\newtheorem{warning}[theorem]{Caution}
\DeclareMathOperator*{\esssup}{ess\,sup}
\newcommand{\E}{\mathbb E}
\newcommand{\Pp}{\mathbb P}
\newcommand{\calH}{\mathcal H}
\newcommand{\calS}{\mathcal S}
\newcommand{\sig}{\sigma}
\newcommand{\eps}{\varepsilon}
\newcommand{\dd}{\mathrm d}
\newcommand{\norm}[1]{\left\lVert#1\right\rVert}
\pagestyle{fancy}
\fancyhf{}
\fancyhead[L]{\small Finite multiplayer stochastic games}
\fancyhead[R]{\small Research note}
\fancyfoot[C]{\thepage}
\setlength{\headheight}{14pt}
\setlist[itemize]{topsep=3pt,itemsep=2pt}
\setlist[enumerate]{topsep=3pt,itemsep=2pt}

\title{On Equilibrium Payoffs in Finite Multiplayer Stochastic Games\\[5pt]
\large Partial Results, Explicit Obstructions, and a Conditional Reduction}
\author{GPT 6 Astra Pro}
\date{8 October 2026}

\begin{document}
\noindent\textbf{Problem ID:} \texttt{C73-1}\quad\textbf{Model:} GPT 6 Astra Pro\quad\textbf{Version:} 1\par
\noindent\textbf{Status:} Unresolved for the complete catalogue problem; this manuscript records partial results and research attempts.\par
\maketitle
\begin{abstract}
The existence of uniform, limiting-average, or liminf-average Nash equilibrium
payoffs in arbitrary finite multiplayer stochastic games remains an open problem.
This note collects and critically audits several attempts to address it. We prove
a sufficient condition based on a state-independent vanishing-discount equilibrium
payoff; exhibit elementary counterexamples to a proposed additive-eigenpair
criterion and to naive concatenation of exact finite-horizon equilibria; and
formulate a history-dependent finite-deviation criterion. Using a finite-prefix
deviator-identification theorem, we prove a conditional conversion from
approximate acceptability and small finite-deviation regret to approximate
liminf-average Nash equilibrium. For each fixed finite deviation horizon we
establish the two required properties simultaneously. The missing step is to
produce one profile satisfying the deviation inequalities for \emph{every}
finite horizon. An example shows why pointwise compactness of strategies cannot
supply this step automatically. No general existence proof or counterexample is
claimed. Statements proved in this note are distinguished from unproved
candidate principles and unsuccessful approaches.
\end{abstract}
\noindent\textbf{Status.} \emph{Research-progress note; not a solution of the general open problem.}
\tableofcontents

\section{The question and precise payoff conventions}
\label{sec:model}
A finite stochastic game is a tuple
\[
\Gamma=(I,S,(A_i(s))_{i\in I,s\in S},q,(g_i)_{i\in I}),
\]
where $I=\{1,\ldots,m\}$ and $S$ are nonempty finite sets, each action set
$A_i(s)$ is finite and nonempty, $q(\cdot\mid s,a)$ is a probability
distribution on $S$, and $g_i(s,a)\in\mathbb R$ is a bounded stage payoff.
At stage $t$ the current state $s_t$ is publicly observed, the players
independently randomize their actions $a_{i,t}$ conditional on the publicly
observed history, receive $g_i(s_t,a_t)$, and the next state is sampled from
$q(\cdot\mid s_t,a_t)$. Actions and states are publicly observed after the
stage. Unless stated otherwise, strategies are arbitrary behavioral strategies
with perfect recall and no external mediator.

Write
\[
\bar g_{i,n}=\frac1n\sum_{t=1}^n g_i(s_t,a_t),\qquad
f_i^-(\omega)=\liminf_{n\to\infty}\bar g_{i,n}(\omega),\qquad
f_i^+(\omega)=\limsup_{n\to\infty}\bar g_{i,n}(\omega).
\]
We use the expected liminf payoff
$u_i^-(s,\sig)=\E_{s,\sig}f_i^-$, with analogous notation $u_i^+$.
One also encounters $\liminf_n\E_{s,\sig}\bar g_{i,n}$ and
$\lim_n\E_{s,\sig}\bar g_{i,n}$ (when the latter exists).
The \emph{expected liminf} and \emph{liminf of expectations} are different:
interchanging expectation and liminf is not justified in general.
A behavioral profile $\sig$ is an $\eps$-equilibrium for $u^-$ from state
$s$ when
\begin{equation}\label{eq:liminfeq}
 u_i^-(s,\tau_i,\sig_{-i})\le u_i^-(s,\sig)+\eps
 \quad\text{for every player }i\text{ and every behavioral }\tau_i.
\end{equation}
A \emph{liminf-average equilibrium payoff} is a vector $c$ such that for
every $\eps>0$ there is an $\eps$-equilibrium (from the specified initial
state) with payoff within $\eps$ of $c$. This is an \emph{approximate}
equilibrium-payoff concept, not an assertion that a zero-error equilibrium
exists.

A profile is a \emph{uniform $\eps$-equilibrium} when, for every sufficiently
large finite horizon $n$, its expected $n$-stage average admits no unilateral
gain above $\eps$, and its payoff is also nearly constant across all sufficiently
large horizons; one may additionally require the analogous inequality for every
sufficiently small normalized discount factor $\lambda$, where
\[
\gamma_i^\lambda(s,\sig)
 =\E_{s,\sig}\sum_{t=1}^\infty
 \lambda(1-\lambda)^{t-1}g_i(s_t,a_t).
\]
A \emph{uniform equilibrium payoff} $c$ is realized asymptotically by
uniform approximate equilibria with their finite-horizon and discounted payoffs
converging to $c$. The formulations above are made precise in the statement
and proof of Theorem~\ref{thm:flat}. We distinguish these concepts rather than
silently identifying them.

\begin{question}[General multiplayer existence problem]\label{q:main}
For every finite stochastic game, and every initial state, does there exist
(a) a uniform equilibrium payoff; or at least (b) a limiting-average equilibrium
payoff; or at least (c) a liminf-average equilibrium payoff in the sense of
\eqref{eq:liminfeq}?
\end{question}
Two-player finite stochastic games have uniform equilibrium payoffs by work of
Vieille; corresponding assertions for arbitrary finite numbers of players are
not established by the references reviewed here
\cite{vieille2002,fleschsolan2023,solanvieille2025,aps2026}.

\section{Established background and a map of the attempts}
For each $\lambda\in(0,1]$, finite discounted stochastic games have stationary
Nash equilibria \cite{fink1964}. Their payoff vectors generally depend on the
starting state, and a discounted equilibrium need not remain approximately
optimal under a limiting-average criterion. In the general-payoff framework,
Flesch and Solan prove the existence of \emph{minmax-acceptable} profiles in
all subgames and of approximate equilibria in at least one subgame, not
necessarily at an arbitrarily prescribed initial state \cite{fleschsolan2023}.
Deviator identification for nearly certain target events is available by
Alon, Gunby, He, Shmaya, and Solan \cite{alon2024}.

\begin{center}
\renewcommand{\arraystretch}{1.18}
\begin{tabular}{@{}p{0.34\linewidth}p{0.57\linewidth}@{}}
\toprule
\textbf{Approach} & \textbf{What survives the audit}\\
\midrule
State-flat discounted branch & A genuine sufficient condition for uniform and liminf-average equilibrium payoffs (Theorem~\ref{thm:flat}).\\
State-independent additive eigenpair & False already for unequal absorbing states (Proposition~\ref{prop:AEfalse}).\\
Repeat exact finite-stage equilibria & False as a general construction: the Big Match supplies a deviation with gain $1/2$ (Proposition~\ref{prop:bigmatch}).\\
Predetermined discount or stopping schedule & Cannot alone secure the Big Match value; observed history matters (Proposition~\ref{prop:hazard}).\\
History-dependent continuation supermartingale & Gives a conditional conversion to liminf-average equilibrium (Theorem~\ref{thm:conversion}).\\
Backward induction for a fixed deviation cutoff & Works for each $N$ separately (Theorem~\ref{thm:prefix}); simultaneous $N$ remains open.\\
Compactness of behavioral strategies & Insufficient: actual tail payoffs may jump under pointwise convergence (Example~\ref{ex:compact}).\\
\bottomrule
\end{tabular}
\end{center}

\section{A valid conditional vanishing-discount theorem}
\label{sec:flat}
Choose a semialgebraic branch $\lambda\mapsto(x^\lambda,v^\lambda)$ of
stationary discounted Nash equilibria for all sufficiently small $\lambda>0$.
Such a branch can be selected because the finite-dimensional discounted
stationary equilibrium correspondence has a semialgebraic graph; a stationary
equilibrium exists for every positive $\lambda$. At finitely many other
discount parameters, choose arbitrary stationary equilibria.
Assume $|g_i|\le M$.

\begin{theorem}[State-flat sufficient condition]\label{thm:flat}
Suppose there are constants $c_i$ such that
\begin{equation}\label{eq:flat}
 \lim_{\lambda\downarrow0}v_i^\lambda(s)=c_i
 \quad\text{for every }i\in I,\ s\in S.
\end{equation}
Thus the limit is independent of the starting state for each player.
Then $c=(c_i)_{i\in I}$ is a uniform equilibrium payoff. The same profiles
also yield equilibrium payoffs $c$ for the expected liminf-average criterion
and for the corresponding limsup-average criterion and the standard
finite-horizon limit-payoff formulations.
\end{theorem}
\begin{proof}
For player $i$, state $s$, and pure action $b_i$, optimality against
$x_{-i}^\lambda$ implies
\begin{equation}\label{eq:bellmandisc}
 v_i^\lambda(s)\ge\lambda
 g_i(s,b_i,x_{-i}^\lambda(s))
 +(1-\lambda)Q_{b_i,x_{-i}^\lambda(s)}v_i^\lambda.
\end{equation}
Here $Q_{b_i,x_{-i}}v=\sum_{s'}q(s'\mid s,b_i,x_{-i})v(s')$.
Equality holds on averaging over $x_i^\lambda(s)$.

With $r\in\mathbb N$, set $v_{i,r}=v_i^{1/r}$ and $x_r=x^{1/r}$.
Multiplying \eqref{eq:bellmandisc} by $r$ yields
\begin{equation}\label{eq:bellmanr}
 r v_{i,r}(s)\ge g_i(s,b_i,x_{-i,r}(s))
 +(r-1)Q_{b_i,x_{-i,r}(s)}v_{i,r}.
\end{equation}
Let $\pi^N$ prescribe $x_r$ when $r$ stages remain in a block of
length $N$. Define
\[
 \Delta_r=\max_{i,s}|v_{i,r-1}(s)-v_{i,r}(s)|,\qquad
 E_N=\sum_{r=2}^N(r-1)\Delta_r.
\]
Induction on $r$, using \eqref{eq:bellmanr}, gives for all starting states
and every unilateral deviation during the block
\begin{equation}\label{eq:blocks}
 \E_{\tau_i,\pi^N_{-i},s}\sum_{t=1}^N g_i
 \le N v_{i,N}(s)+E_N,
\end{equation}
and, using equality under the prescribed mixed actions,
\begin{equation}\label{eq:blockbaseline}
 \left|\E_{\pi^N,s}\sum_{t=1}^N g_i-Nv_{i,N}(s)\right|
 \le E_N.
\end{equation}
The induction works for arbitrary within-block history-dependent deviations:
a best-response continuation can be optimized at each successor history,
while $\Delta_r$ bounds the error from replacing $v_{i,r-1}$ by $v_{i,r}$.

A bounded one-variable semialgebraic function with a finite limit has a
Puiseux expansion near zero. In particular, for some $\alpha>0$ and $C<\infty$,
\[
 \left\|\frac{\dd v_i^\lambda}{\dd\lambda}\right\|_\infty
 \le C\lambda^{\alpha-1}
\]
for sufficiently small $\lambda>0$. Therefore
$\Delta_r=O(r^{-1-\alpha})$ and
\begin{equation}\label{eq:error}
 E_N=O\left(\sum_{r=2}^Nr^{-\alpha}\right)+O(1)=o(N).
\end{equation}
Under \eqref{eq:flat}, the errors
\[
 d_N=\max_{i,s}|v_{i,N}(s)-c_i|+E_N/N
\]
converge to zero. Consequently every $N$-block, from every initial state,
has equilibrium expected average within $d_N$ of $c_i$, and a deviator's
expected average is at most $c_i+d_N$.

Repeat $\pi^N$ periodically, restarting only the \emph{countdown} after each
block; the physical state is \emph{not} reset. Call this profile $\sig^N$.
Conditioning at the start of each block and summing the bounds gives, for
all $n$,
\begin{align}
 \E_{s,\sig^N}\bar g_{i,n}&\ge c_i-d_N-2MN/n,\label{eq:uniformlower}\\
 \E_{s,\tau_i,\sig^N_{-i}}\bar g_{i,n}
 &\le c_i+d_N+2MN/n.\label{eq:uniformupper}
\end{align}
The same argument gives an upper bound on the prescribed payoff. Since
$N$ is fixed while $n\to\infty$, this proves the long finite-horizon part
of uniform equilibrium.

For completeness, the discounted criterion follows from the same block
bounds. Let $\theta\downarrow0$ be an \emph{evaluation} discount factor,
distinct from the within-block parameters $1/r$. The normalized discount
weights of successive complete blocks satisfy
\[
 \frac{\theta\sum_{t=kN+1}^{(k+1)N}(1-\theta)^{t-1}g_t}
 {\theta\sum_{t=kN+1}^{(k+1)N}(1-\theta)^{t-1}}
 =\frac1N\sum_{t=kN+1}^{(k+1)N}g_t+O(MN\theta),
\]
uniformly in the block index and realization; the geometric weights within
one block have relative variation $O(N\theta)$. The block inequalities
then imply that discounted prescribed payoffs are within
$d_N+O(MN\theta)$ of $c_i$ and deviations yield at most
$c_i+d_N+O(MN\theta)$. First take $N$ large, then $\theta$ small.

For the pathwise liminf/limsup criteria, let $Y_{i,k}$ be the total reward
in block $k$. Under either prescribed play or one fixed arbitrary unilateral
deviation, the bounded martingale differences
\[
 Y_{i,k}-\E[Y_{i,k}\mid\calH_{kN}]
\]
have averages converging almost surely to zero. Under prescribed play,
all conditional block expectations are within $Nd_N$ of $Nc_i$;
under a unilateral deviation, they are at most $N(c_i+d_N)$.
Thus almost surely under $\sig^N$,
\[
 c_i-d_N\le\liminf_n\bar g_{i,n}
 \le\limsup_n\bar g_{i,n}\le c_i+d_N,
\]
and under any $\tau_i$,
$\limsup_n\bar g_{i,n}\le c_i+d_N$ almost surely.
Bounded convergence inequalities then show that $\sig^N$ is a
$2d_N$-equilibrium for both $\E\liminf$ and $\E\limsup$, with payoffs
within $d_N$ of $c$. For formulations using $\liminf_n\E\bar g_{i,n}$ or
$\limsup_n\E\bar g_{i,n}$, the same conditional-block estimates
apply directly. For the prescribed periodic finite-memory profile,
the ordinary limit $\lim_n\E\bar g_{i,n}$ exists by finite-state
Markov-chain ergodic theory. An arbitrary deviating strategy need not
have a well-defined ordinary limit of expected average payoffs; no
unqualified assertion is made here for a criterion that is undefined
under such deviations. This completes the claimed approximate-payoff
conclusions.
\end{proof}

\begin{remark}
The state-flat hypothesis is essential to this proof. Without it,
\eqref{eq:blocks} bounds the payoff in block $k$ by
$N v_{i,N}(S_{kN+1})+E_N$. A deviator can change the next block's
starting-state distribution. The state-dependent benchmarks need not
telescope or be bounded by one common target. Thus Theorem~\ref{thm:flat}
is a sufficient result, not a proof of Question~\ref{q:main}.
\end{remark}

\section{Two false generalizations of the discounted argument}
\subsection{A state-independent additive eigenpair need not exist}
An earlier attempt asked for a constant $c\in\mathbb R^m$, a potential
$\Phi:S\to\mathbb R^m$, a length $N$, and Nash equilibrium payoffs
$y(s)$ in the $N$-stage game with terminal payment $\Phi(S_{N+1})$ such that
\begin{equation}\label{eq:AE}
 \max_{s\in S}\norm{y(s)-\Phi(s)-Nc}_\infty\le N\eta.
\tag{AE}
\end{equation}
If valid, the terminal potentials would telescope across blocks. But the
assertion that \eqref{eq:AE} holds for every finite stochastic game and
arbitrarily small $\eta$ is false.

\begin{proposition}\label{prop:AEfalse}
For $\eta<\tfrac12$, there is a finite stochastic game for which no $N$,
$c$, or $\Phi$ satisfies \eqref{eq:AE}.
\end{proposition}
\begin{proof}
Take two absorbing states $z_0,z_1$, no meaningful choices, and reward $0$
per stage at $z_0$ and reward $1$ per stage at $z_1$ for every player.
From $z_j$ the only terminal-reward payoff is
$y_i(z_j)=Nj+\Phi_i(z_j)$. Condition \eqref{eq:AE} would imply
$|c_i|\le\eta$ and $|1-c_i|\le\eta$, a contradiction.
\end{proof}
\begin{remark}
This is \emph{not} a counterexample to equilibrium existence: every
profile is an equilibrium. It only rules out a state-independent
additive-eigenpair principle of the stated form.
\end{remark}

\subsection{Exact countdown blocks fail in the Big Match}
\label{sec:bigmatch}
Consider the two-player constant-sum Big Match, with player 1's reward
shown below and player 2's reward equal to $1-g_1$. A star denotes
permanent absorption at the indicated payoff:
\begin{equation}\label{eq:BM}
\begin{array}{c|cc}
 & L&R\\\hline
 T&1^{*}&0^{*}\\
 B&0&1
\end{array}.
\end{equation}
At $r$ remaining stages, if the value of the continuation with $r-1$
stages is $(r-1)/2$, the matrix of \emph{total} continuation rewards is
\[
\begin{array}{c|cc}
 &L&R\\\hline
 T&r&0\\
 B&(r-1)/2&(r+1)/2
\end{array}.
\]
Player 1 plays $T$ with probability $p_r=1/(r+1)$ and player 2 mixes
$L,R$ equally. Direct indifference checks and induction show that this is
an \emph{exact} equilibrium of the $N$-stage game with value $N/2$.
It agrees with the countdown $x^{1/r}$ of stationary discounted equilibria:
\[
 v^\lambda=\tfrac12\quad\text{at the nonabsorbing state},
 \qquad p^\lambda=\frac{\lambda}{1+\lambda}.
\]
The state-flat hypothesis fails: the two absorbing states have discounted
payoffs $0$ and $1$.

\begin{proposition}[Exact blocks, fixed profitable deviation]
\label{prop:bigmatch}
Repeat the exact $N$-stage equilibrium countdown block forever, restarting
the countdown after each $N$ stages but not resetting the physical state.
For every $N$, player 2 has a unilateral deviation increasing expected
liminf-average payoff by $1/2$.
\end{proposition}
\begin{proof}
Conditional on entering a block without prior absorption, the probability
of surviving its $N$ stages is
\begin{equation}\label{eq:survival}
 \prod_{r=1}^N\left(1-\frac1{r+1}\right)
 =\prod_{r=1}^N\frac r{r+1}=\frac1{N+1}.
\end{equation}
Thus absorption occurs almost surely during the infinite concatenation.
Under prescribed play, the absorbing column is fair, so both players'
expected limiting average payoffs equal $1/2$. If player 2 plays $R$
forever, the stopping hazard is unchanged, but absorption always yields
player 1 reward $0$ and player 2 reward $1$. The deviation gain of player 2
is exactly $1/2$ for every $N$.
\end{proof}
A quantitative finite-horizon version makes the failure especially clear.
Against the deviation $R$ forever, player 1's expected total reward in
each reached block equals $N/2$. By \eqref{eq:survival}, the expected total
reward over \emph{all} time is $(N/2)\sum_{k\ge0}(N+1)^{-k}=(N+1)/2$.
Therefore $\E\bar g_{1,n}\le(N+1)/(2n)$, whereas the prescribed
profile yields expected average $1/2$. Exact block Nash optimality does
not prevent deviations with effects that persist after the nominal block.

\begin{proposition}[Predetermined stopping hazards are insufficient]
\label{prop:hazard}
Suppose player 1's stopping probability at date $t$, conditional on
nonabsorption, is a fixed number $p_t$, independent of observed play.
Then player 2 can force player 1's \emph{expected liminf-average reward}
to be arbitrarily close to $0$.
\end{proposition}
\begin{proof}
If $\sum_t p_t=\infty$, playing $R$ forever causes absorption at payoff
$0$ almost surely. If $\sum_t p_t<\infty$, choose $K$ with
$\sum_{t>K}p_t<\delta$; player 2 plays $R$ through $K$ and $L$
thereafter. Absorption before $K$ yields $0$, while positive limiting
payoff is possible only if player 1 stops after $K$. The probability of
that event is at most $\sum_{t>K}p_t<\delta$. Since payoffs lie in
$[0,1]$, player 1's expected liminf-average payoff is at most $\delta$.
\end{proof}
The Big Match itself has approximate long-run equilibria; the propositions
refute \emph{time-only implementations}, not its equilibrium existence
\cite{bigmatchspace}.

\subsection{Discounted equilibrium limits need not be uniform payoffs}
Even without the block-concatenation error, a direct assertion that every
vanishing-discount Nash payoff limit is uniform is false: Sorin's example,
discussed by Renault and Ziliotto, has a nonempty discounted equilibrium
payoff limit set disjoint from the set of uniform equilibrium payoffs
\cite{renaultziliotto}. This is a distinct reason why a universal
vanishing-discount proof needs additional hypotheses or a different target.

\section{A history-dependent, finite-deviation conversion theorem}
\label{sec:conversion}
The previous failures suggest avoiding one state-independent payoff target
at \emph{every} intermediate date. We now let the continuation benchmark
be an actual payoff conditional on the entire finite public history.
Normalize rewards in this and the remaining sections to $[0,1]$.

For a profile $\sig$, history $h$, and player $i$, put
\begin{equation}\label{eq:V}
 V_i^\sig(h)=\E_{h,\sig}f_i^-.
\end{equation}
Let
\begin{equation}\label{eq:minmax}
 v_i(s)=\inf_{\rho_{-i}}\sup_{\tau_i}
 \E_{s,\tau_i,\rho_{-i}}f_i^-
\end{equation}
be the individual \emph{minmax} level. Opponents in $\rho_{-i}$ use an
ordinary joint \emph{profile} of independently randomized behavioral
strategies, rather than a mediator. Since $f_i^-$ is shift-invariant, the
minmax level after a finite history depends only on the current state.
By the definition of the infimum, for every $\zeta>0$ there is a
joint-profile punishment against $i$ yielding at most $v_i(s)+\zeta$
against any strategy of $i$.

Let $\tau_i\star_N\sig_i$ be a strategy that may differ from $\sig_i$
through stage $N$ and equals $\sig_i$ at all subsequent histories. Define
\begin{align}
 A(\sig)&=\sup_{i,h}[v_i(s_h)-V_i^\sig(h)]_+,\label{eq:A}\\
 R_N(\sig)&=\max_i\sup_{\tau_i}\Bigl[
  u_i^-(s_0,\tau_i\star_N\sig_i,\sig_{-i})
   -u_i^-(s_0,\sig)\Bigr],\label{eq:RN}\\
 R_{\mathrm{fin}}(\sig)&=\sup_{N\ge1}R_N(\sig).\label{eq:Rfin}
\end{align}
These $R_N$ use \emph{infinite-horizon tail payoffs}; only the length of
the deviation is bounded. This is not an $N$-stage expected-payoff test.

\begin{theorem}[Finite-deviation conversion; conditional on a starting
profile]\label{thm:conversion}
Suppose $A(\sig)\le\alpha$ and $R_{\mathrm{fin}}(\sig)\le\beta$.
For every $\rho,\delta\in(0,1)$ there is an ordinary behavioral profile
$\widehat\sig$ such that, from $s_0$,
\begin{align}
 |u_i^-(\widehat\sig)-u_i^-(\sig)|&\le\delta
 &&\text{for every }i,\label{eq:convpay}\\
 \sup_{\tau_i}\bigl[
 u_i^-(\tau_i,\widehat\sig_{-i})-u_i^-(\widehat\sig)
 \bigr]
 &\le\alpha+\beta+\rho+2\sqrt{m\delta}+\delta
 &&\text{for every }i.\label{eq:conveps}
\end{align}
Here $m=|I|$; the estimate uses the finite-prefix identification
form of \cite[Theorem~2.6 and Section~4]{alon2024}, with stochastic
transitions treated as prescribed moves of Nature in an alternating
perfect-information expansion. In particular, the conclusion does
\emph{not} add an actual mediator to the game.
\end{theorem}
\begin{proof}
\emph{Step 1: a continuation supermartingale.}
For a finite history $h$ define
\begin{equation}\label{eq:W}
 W_i(h)=\sup_{N\ge |h|}\sup_{\tau_i\text{ through }N}
 \E_{h,\tau_i,\sig_{-i}}\!
 \left[V_i^\sig(H_N)\right].
\end{equation}
The player returns to $\sig_i$ after $N$. (The indexing of $H_N$ may
be shifted by one without changing the argument.) Then
\begin{equation}\label{eq:Wbound}
 0\le V_i^\sig(h)\le W_i(h)\le1,
 \qquad W_i(s_0)\le u_i^-(\sig)+\beta.
\end{equation}
Concatenating a current pure deviation with almost optimal bounded-time
deviations from the finitely many successor histories, using one common
terminal date, shows
\begin{equation}\label{eq:Wsuper}
 \E\left[W_i(H_{t+1})\mid H_t=h,a_i,\sig_{-i}(h)\right]
 \le W_i(h).
\end{equation}
Thus, when the opponents use $\sig_{-i}$, $W_i(H_t)$ is a bounded
supermartingale under \emph{every} unilateral deviation. The common-date
argument is legitimate because at any stage there are finitely many
successor public histories.

\emph{Step 2: a compact full-payoff-realization set.}
Under $\sig$, the bounded martingale
$V_i^\sig(H_t)=\E_\sig[f_i^-\mid\calH_t]$ converges almost surely to
$f_i^-$. Accordingly, if
\[
 G=\{\omega: V_i^\sig(H_t(\omega))\to f_i^-(\omega)
 \text{ for every }i\},
\]
then $\Pp_\sig(G)=1$. The infinite-play space is compact metrizable;
regularity of its probability measures gives a compact $K\subset G$
with $\Pp_\sig(K)>1-\delta$.
Let $T$ be the first finite public history whose cylinder has empty
intersection with $K$. Closedness of $K$ implies
\begin{equation}\label{eq:exit}
 \{T=\infty\}=K.
\end{equation}

\emph{Step 3: identify an offending player at a finite exit prefix.}
Apply the closed-target finite-prefix blame theorem
\cite[Section~4]{alon2024} to the prescribed law and target $K$.
The standard reduction of stochastic transitions to an additional
non-strategic Nature agent is performed by splitting each stage into
(i) simultaneous player actions and (ii) Nature's publicly observed
state draw according to $q$. Nature has a prescribed behavioral strategy
and does not deviate. There are $m+1$ agents in the identification
problem; map any attribution to Nature to an arbitrary genuine player.
We obtain a \emph{finite-exit-prefix-measurable} blame assignment
$j(H_T)\in I$ such that, for each genuine sole deviator $i$,
\begin{equation}\label{eq:blame}
 \Pp_{\tau_i,\sig_{-i}}\{T<\infty,\ j(H_T)\ne i\}
 \le 2\sqrt{m\delta}.
\end{equation}
At $T$, modify the continuation: the players other than $j(H_T)$
switch permanently to a $\rho$-minmax punishment of that player from
the current state. Call the resulting profile $\widehat\sig$.
The event $T<\infty$ has prescribed probability at most $\delta$;
therefore boundedness of $f_i^-$ gives \eqref{eq:convpay}.

\emph{Step 4: bound an arbitrary infinite deviation.}
Fix $i$ and an unrestricted $\tau_i$. Before $T$, the opponents use
$\sig_{-i}$; hence $W_i(H_{t\wedge T})$ is a bounded supermartingale,
converging almost surely to some $Z_i$ with
$\E Z_i\le W_i(s_0)$. On $\{T=\infty\}$, the path lies in $K$, so
\begin{equation}\label{eq:onK}
 f_i^-=\lim_tV_i^\sig(H_t)
 \le\lim_tW_i(H_t)=Z_i.
\end{equation}
On $\{T<\infty,j(H_T)=i\}$, the punishment yields conditional expected
payoff at most
\begin{equation}\label{eq:correctblame}
 v_i(s_{H_T})+\rho
 \le V_i^\sig(H_T)+\alpha+\rho
 \le W_i(H_T)+\alpha+\rho.
\end{equation}
On incorrect blame use the payoff bound $1$ and \eqref{eq:blame}.
Because the law of the stopped history is unaffected by the punishment
that follows the exit,
\[
 u_i^-(\tau_i,\widehat\sig_{-i})
 \le\E Z_i+\alpha+\rho+2\sqrt{m\delta}
 \le u_i^-(\sig)+\beta+\alpha+\rho+2\sqrt{m\delta}.
\]
Together with \eqref{eq:convpay}, this proves \eqref{eq:conveps}.
\end{proof}

\begin{remark}[Nature and observability]
The identification step requires public observation of realized actions
and states and a known transition law. It should not be applied verbatim
to games with hidden actions, private states, or an endogenous unobserved
mediator. The alternating Nature construction is a device for applying a
finite-prefix testing theorem; it does not expand the feasible deviations
in the original game.
\end{remark}

\subsection{The Big Match passes the finite-deviation test}
In \eqref{eq:BM}, prescribe player 1 to play $B$ forever and player 2
to choose independent fair $L/R$ forever. At every unabsorbed history,
\[
 V^\sig(h)=(1/2,1/2),
\]
and at absorbing histories $V^\sig$ is the absorbing reward vector.
The minmax value to each player is $1/2$ in the nonabsorbing state and
the fixed terminal reward after absorption; hence $A(\sig)=0$.
Every bounded-duration deviation of player 1 either leads to an absorption
whose \emph{expected} payoff is $1/2$, or returns to continuation payoff
$1/2$. Every bounded-duration deviation of player 2 affects only finitely
many rewards, because player 1 never stops. Thus
\begin{equation}\label{eq:BMcriteria}
 A(\sig)=R_{\mathrm{fin}}(\sig)=0.
\end{equation}
Yet this $\sig$ is not itself an equilibrium: player 2 can play $L$
forever to obtain payoff $1$. Theorem~\ref{thm:conversion} yields
approximate equilibria by testing deviations and assigning punishment.
This is a consistency check on the \emph{conditional} conversion theorem,
not a new solution of the Big Match.

\section{A fixed-cutoff construction in every finite game}
\label{sec:prefix}
The following is the strongest unconditional step from the
finite-deviation approach.

\begin{theorem}[Finite-prefix construction]\label{thm:prefix}
For every finite stochastic game, every $\alpha>0$, and every integer
$N\ge1$, there is a behavioral strategy profile $\sig^{(N)}$ satisfying
\begin{equation}\label{eq:prefix}
 A(\sig^{(N)})\le\alpha,
 \qquad R_N(\sig^{(N)})=0.
\end{equation}
\end{theorem}
\begin{proof}
The minmax-acceptability theorem of Flesch and Solan
\cite[Theorem~4.4]{fleschsolan2023} supplies a profile $\pi$ with
\begin{equation}\label{eq:acceptablepi}
 V_i^\pi(h)\ge v_i(s_h)-\alpha
 \qquad\text{for all players and all finite histories }h.
\end{equation}
Construct a finite extensive-form game for the first $N$ stages, with
terminal payoff $V_i^\pi(h_N)$ at each date-$N$ history; after date $N$
prescribe $\pi$. Choose a mixed Nash equilibrium in every finite
continuation game by backward induction and call the full profile
$\sig^{(N)}$.

A deviation ending by date $N$ cannot improve player $i$'s expected
tail payoff: conditional on the date-$N$ history, the remaining payoff
is precisely $V_i^\pi(h_N)$. Thus $R_N(\sig^{(N)})=0$.

To propagate acceptability backward, use the dynamic-programming identity
for the state-dependent minmax level with shift-invariant payoff:
\begin{equation}\label{eq:minmaxBellman}
 v_i(s)=\min_{x_{-i}\in\prod_{j\ne i}\Delta(A_j(s))}
 \max_{a_i\in A_i(s)}
 \sum_{a_{-i},s'}x_{-i}(a_{-i})q(s'\mid s,a_i,a_{-i})v_i(s').
\end{equation}
For clarity, \eqref{eq:minmaxBellman} uses \emph{product} mixed actions
of the opponents. It follows by decomposing an arbitrary opponent
strategy into its first-stage product mixed action and continuation
profiles and using the shift-invariance of $f_i^-$. The upper bound
concatenates near-minmax opponent continuation profiles after each
observed action-state history; the lower bound responds optimally at
the first stage and then nearly best-responds to the opponents'
continuations. Only finitely many first-stage successors arise, so the
approximation errors can be made arbitrarily small.

If all next-history continuation values $y_i(h')$ satisfy
$y_i(h')\ge v_i(s_{h'})-\alpha$, the current one-shot Nash equilibrium
with opponent product mixed action $x_{-i}$ gives
\[
\begin{split}
 y_i(h)
 &\ge\max_{a_i}\E[y_i(H_{t+1})\mid h,a_i,x_{-i}]\\
 &\ge\max_{a_i}\E[v_i(S_{t+1})\mid h,a_i,x_{-i}]-\alpha\\
 &\ge v_i(s_h)-\alpha,
\end{split}
\]
by \eqref{eq:minmaxBellman}. This inducts backward from
\eqref{eq:acceptablepi}. After date $N$, acceptability again follows
from \eqref{eq:acceptablepi}.
\end{proof}

\section{The remaining quantifier gap and a compactness warning}
\label{sec:gap}
By Theorem~\ref{thm:conversion}, the following \emph{unproved}
selection principle would settle the liminf-average part of
Question~\ref{q:main}:
\begin{question}[Simultaneous finite-deviation criterion]\label{q:criterion}
Does every finite stochastic game, for each initial state $s_0$ and
$\eta>0$, admit a single profile $\sig$ such that
\begin{equation}\label{eq:missing}
 A(\sig)\le\eta,
 \qquad \sup_{N\ge1}R_N(\sig)\le\eta\ ?
\tag{FD}
\end{equation}
\end{question}
If \eqref{eq:missing} held for arbitrarily small $\eta$, the conversion
theorem would give approximate liminf-average equilibria. Compactness of
the \emph{numerical payoff vectors} (not of their strategy realizations)
would then yield a liminf-average equilibrium payoff in the approximate
sense defined in Section~\ref{sec:model}.

Theorem~\ref{thm:prefix} proves
\begin{equation}\label{eq:quantifiers}
 \forall\eta>0\ \forall N\ \exists\sig^{(N)}:
 \quad A(\sig^{(N)})\le\eta,\ R_N(\sig^{(N)})=0,
\end{equation}
whereas \eqref{eq:missing} asks for
\[
 \forall\eta>0\ \exists\sig\ \forall N:
 \quad A(\sig)\le\eta,\ R_N(\sig)\le\eta.
\]
The quantifier interchange is \emph{not} justified.

\begin{proposition}[Tail payoffs need not survive pointwise strategy limits]
\label{ex:compact}
There is a one-player two-state stochastic game and a sequence of exact
equilibrium strategies $\sig^{(N)}$ converging pointwise at every fixed
history to a strategy $\sig^\infty$, such that
\[
 u^-(\sig^{(N)})=1\quad\forall N,
 \qquad u^-(\sig^\infty)=0.
\]
\end{proposition}
\begin{proof}
At nonabsorbing state $s$, \emph{wait} yields reward $0$ and remains in
$s$, while \emph{stop} moves to an absorbing state of permanent reward
$1$. Let $\sig^{(N)}$ wait through date $N-1$ and stop at date $N$;
after any off-path history that remains unabsorbed at or after that date,
stop immediately. Every such strategy is optimal and has tail payoff
$1$ from every history. On each fixed finite history, however, as
$N\to\infty$ its prescribed action is eventually \emph{wait}. The
pointwise limit waits forever and gets tail payoff $0$.
\end{proof}
This example does not disprove \eqref{eq:missing}: choosing \emph{stop}
immediately satisfies it. It proves only that an arbitrary strategy-limit
argument fails. The key requirement is to construct a profile for which
\emph{actual realized tail payoffs} and \emph{all finite-deviation
incentive inequalities} survive simultaneously.

\section{Research directions and boundaries of the present note}
\begin{enumerate}[label=\textbf{\arabic*.},leftmargin=*]
\item \textbf{History-dependent selection.} Find a direct construction or a
compactness theorem tailored to actual continuation payoffs
$V_i^\sig(h)$ and the entire sequence $(R_N)_{N\ge1}$.
A mere subsequence of behavior profiles does not suffice by
Proposition~\ref{ex:compact}.
\item \textbf{Tests before profitable steering.} Develop a deviation test
that limits the probability of reaching higher individual minmax levels
before the player receives effective punishment. This is the central
state-dependent continuation obstruction highlighted in
\cite{fleschsolan2023}.
\item \textbf{Do not conflate solution concepts.} Correlated equilibrium,
existence from \emph{some} initial state, and existence at \emph{every}
initial state are different assertions. A general correlated-existence
theorem does not supply the desired ordinary Nash profile.
\item \textbf{Restricted positive results.} Specialized approaches for
absorbing and quitting games are informative but do not answer the
arbitrary finite-state multiplayer question. Recent examples include
\cite{solanvieille2025,aps2026}.
\end{enumerate}

\begin{remark}[What is and is not claimed]
Theorem~\ref{thm:flat}, Propositions~\ref{prop:AEfalse},
\ref{prop:bigmatch}, \ref{prop:hazard},
Theorem~\ref{thm:conversion}, and Theorem~\ref{thm:prefix}
are mathematical statements proved here, the conversion theorem relying
on the published deviator-identification and acceptability results
explicitly cited. The condition \eqref{eq:missing} remains \emph{unproved}.
Neither the state-flat result nor the finite-cutoff construction, taken
alone or together, resolves Question~\ref{q:main}. No counterexample to
the open existence statement is exhibited here.
\end{remark}

\section*{References and source notes}
\addcontentsline{toc}{section}{References and source notes}
\begin{thebibliography}{99}
\bibitem{fink1964}
A. M. Fink, \emph{Equilibrium in a stochastic $n$-person game},
Journal of Science of the Hiroshima University, Series A-I Mathematics
\textbf{28} (1964), 89--93.

\bibitem{vieille2002}
N. Vieille, \emph{Stochastic games: Recent results},
Handbook of Game Theory with Economic Applications, Vol. 3 (2002),
1833--1850. DOI: \href{https://doi.org/10.1016/S1574-0005(02)03011-4}{10.1016/S1574-0005(02)03011-4}.

\bibitem{fleschsolan2023}
J. Flesch and E. Solan, \emph{Stochastic Games with General Payoff Functions},
Mathematics of Operations Research \textbf{49} (2024), 1349--1371,
published online 2023.
DOI: \href{https://doi.org/10.1287/moor.2023.1385}{10.1287/moor.2023.1385}.
Preprint: \url{https://arxiv.org/abs/2208.12096}.

\bibitem{alon2024}
N. Alon, B. Gunby, X. He, E. Shmaya, and E. Solan,
\emph{Identifying the Deviator}, arXiv:2203.03744, revised 2024.
\url{https://arxiv.org/abs/2203.03744}.
In particular, Theorem~2.6 and the finite-prefix construction in Section~4.

\bibitem{bigmatchspace}
K. A. Hansen, R. Ibsen-Jensen, and M. Kouck\'{y},
\emph{The Big Match in Small Space}, arXiv:1604.07634.
\url{https://arxiv.org/abs/1604.07634}.

\bibitem{renaultziliotto}
J. Renault and B. Ziliotto,
\emph{Hidden Stochastic Games and Limit Equilibrium Payoffs},
arXiv:1407.3028.
\url{https://arxiv.org/abs/1407.3028}.
The introduction discusses Sorin's example separating discounted-equilibrium
payoff limits from uniform equilibrium payoffs in an ordinary stochastic game;
the paper also studies hidden stochastic games, which must not be confused with
the finite perfect-monitoring model in Question~\ref{q:main}.

\bibitem{solanvieille2025}
E. Solan and N. Vieille,
\emph{Undiscounted Equilibrium in Positive Recursive Absorbing Games with
Non-Rectangular Absorption Structure}, arXiv:2512.04306 (2025).
\url{https://arxiv.org/abs/2512.04306}.

\bibitem{aps2026}
G. Ashkenazi-Golan, I. Krasikov, C. Rainer, and E. Solan,
\emph{The APS approach for undiscounted quitting games},
International Journal of Game Theory \textbf{55} (2026), article 19.
DOI: \href{https://doi.org/10.1007/s00182-026-00982-6}{10.1007/s00182-026-00982-6}.
\end{thebibliography}
\section{Completion Estimate}
25\%

\end{document}
